\documentclass[11pt]{article}
\usepackage[margin=1in]{geometry}
\usepackage[T1]{fontenc}

\usepackage{amsmath,amssymb,amsthm}
\usepackage{mathtools}
\usepackage{booktabs}
\usepackage{multirow}
\usepackage{graphicx}
\usepackage{xcolor}
\usepackage{algorithm}
\usepackage{algpseudocode}
\usepackage[numbers,sort&compress]{natbib}
\usepackage[hidelinks]{hyperref}
\usepackage{siunitx}
\sisetup{group-separator={,},group-minimum-digits=4}

\newtheorem{theorem}{Theorem}
\newtheorem{lemma}[theorem]{Lemma}
\newtheorem{proposition}[theorem]{Proposition}
\newtheorem{corollary}[theorem]{Corollary}
\theoremstyle{definition}
\newtheorem{definition}[theorem]{Definition}
\newtheorem{remark}[theorem]{Remark}

\newcommand{\OPT}{\mathrm{OPT}}
\newcommand{\cost}{\mathrm{cost}}
\newcommand{\LB}{\mathrm{LB}}
\newcommand{\Ep}{E^{+}}
\newcommand{\Nt}{N_{2}}
\newcommand{\R}{\mathbb{R}}

\title{Certified Correlation Clustering at Scale:\\
Distance-Two Relaxations, Anytime Dual Bounds,\\ and Exact Partition Crossover}

\author{Quang-Vinh Dang\\
British University Vietnam, Hung Yen, Vietnam\\
\texttt{vinh.dq4@buv.edu.vn} \quad ORCID: 0000-0002-3877-8024}

\date{}

\begin{document}
\maketitle

\begin{abstract}
Correlation clustering on complete signed graphs (equivalently, cluster
editing) has seen rapid progress on approximation ratios.  LP and SDP
hierarchies now give ratios below $1.5$, and a combinatorial algorithm
achieves $2-2/13$ in near-linear time.  In practice, however, the solutions of
engineered heuristics come with no quality guarantee on large graphs.  Exact
solvers stop at a few hundred vertices, and the metric LP has $\Theta(n^3)$
constraints.  The strongest combinatorial lower bounds, the packing bounds of
branch-and-bound solvers, store the instance as a dense matrix, and the
bounds that do scale certify only a factor of about two.  We close this gap
from both sides.
(i)~Optimal clusters have diameter at most two.  We use this to restrict the
metric LP to the $O(\sum_v \deg(v)^2)$ pairs at distance at most two, without
losing validity or the guarantees of LP-rounding algorithms, and we
strengthen it with star and local subgraph inequalities.
(ii)~We compute an anytime, always-valid lower bound by dual
block-coordinate ascent.  It is warm-started by sparse star packings improved
with ruin-and-recreate local search, and every block is an exact LP.
(iii)~We prove an exact decomposition of the optimality gap into non-negative
per-pair and per-constraint terms.  It yields a local optimality certificate
and guides a large-neighbourhood search with exact sub-MIPs.
Together with the first implementation of the iterated flipping local search
of Cohen-Addad et al., these ingredients form \emph{CertiFlip}.  CertiFlip
has an a-priori guarantee ($3$ in expectation, and $2.06$ whenever the
relaxation is solved to optimality) and outputs a certified bound on its own
approximation ratio.  On the PACE~2021 exact track it improves the best published
bounds of three open instances, and its bound proves optimality on more
instances than the root bounds of the state-of-the-art branch-and-bound.  On SNAP graphs with up to $10^6$
edges, its bounds certify the best known solutions within $0.2\%$ to $9.2\%$
of optimal, where previously scalable bounds certify only factors of $1.17$ to
$1.81$.
(iv)~We show that two clusterings can be recombined exactly: in every block
of the overlap of their clusters one may take either parent, and the best
choice is never worse than either parent.  With the contraction of critical
cliques, annealing with region-wise acceptance and localized annealing with
exact rollback, this gives a memetic solver, \emph{PXMem}, with the same
guarantee of $3$ in expectation.  In 110 head-to-head runs of 600\,s against
the PACE~2021 winner KaPoCE on the same machine, PXMem is significantly better
on four of eleven graphs, significantly worse on one (by $0.004\%$), and
equivalent on six.  The same crossover combines parallel runs and improves
the solutions of KaPoCE itself.
(v)~The combinatorial statements behind both solvers, including the move
formulas used by the implementation, the lower-bound certificates and the
approximation guarantee, are machine-checked in Lean~4.
\end{abstract}

\section{Introduction}

In \emph{correlation clustering}~\citep{BansalBC04} we are given a complete
graph whose pairs are labelled \emph{similar} ($+$) or \emph{dissimilar}
($-$) and look for a partition of the vertices minimising the number of
disagreements: $+$ pairs between clusters plus $-$ pairs inside clusters.  The
number of clusters is not prescribed.  Equivalently, the input is the positive
graph $G^+=(V,\Ep)$ and the task is to add or delete the fewest edges so that
$G^+$ becomes a disjoint union of cliques (\emph{cluster editing}).  The
problem is APX-hard~\citep{CharikarGW05} and arises in entity resolution,
community detection and the aggregation of inconsistent information.

\paragraph{The state of the art is split.}  On the theory side, the
approximation ratio for complete graphs dropped from $3$ for the Pivot
algorithm~\citep{AilonCN08} to $2.06$ by metric-LP
rounding~\citep{ChawlaMSY15}, below $2$ with Sherali--Adams
hierarchies~\citep{CohenAddadLN22}, to $1.73$ with
preclustering~\citep{CohenAddadLLN23} and to $1.485$ with the cluster
LP~\citep{CaoCLLNV24}.  A preprint from July~2026 reports a deterministic
$1.3865$-approximation and an upper bound of $1.3865$ on the integrality gap
of the cluster LP~\citep{GarciaSorianoS26}.  The best hardness is
$24/23-\varepsilon$ unless $\mathrm{P}=\mathrm{BPP}$~\citep{CaoCLLNV24}.
Solving the cluster LP takes $2^{\mathrm{poly}(1/\varepsilon)}$ factors even in
the sublinear algorithm of \citet{CaoCLLLNTVYZ25}.  The best
\emph{combinatorial} ratio is $2-2/13$, achieved by a local search with
``flips'' in near-linear time~\citep{CohenAddadLPTYZ24}.  As far as we could
determine, none of the sub-$3$ near-linear algorithms has a public
implementation or has been evaluated experimentally.

On the practical side, engineered heuristics are very strong.  The winner of the
heuristic track of the PACE~2021 challenge, KaPoCE~\citep{BlasiusEtAl21kapoceHeur},
and multilevel memetic solvers~\citep{hausberger2025scc} find excellent
solutions.  However, nobody can say \emph{how} good these solutions are on
large graphs.  Exact branch-and-bound solves instances with a few hundred
vertices~\citep{BlasiusEtAl22sea}.  Metric-LP solvers based on projection
methods reach about $10^4$ vertices~\citep{veldt2019metric,ruggles2020parallel}.
The branch-and-bound algorithm of \citet{BlasiusEtAl22sea} uses strong
packing lower bounds, but its implementation stores the instance as a dense
matrix.  The combinatorial bounds that scale to large graphs certify
solutions only within a factor of about two~\citep{veldt2022stc}.  The PACE report~\citep{KellerhalsKNZ21pace}
lists 27 exact-track instances with up to 620 vertices whose optimum is
still unknown.

\paragraph{This paper.}  We ask for a \emph{practical algorithm with a
provable guarantee that also certifies its own quality on the instance at
hand}, in the sense of certifying algorithms~\citep{McConnellMNS11}.  Our
contributions are as follows.

\begin{enumerate}
\item \textbf{Distance-two relaxations} (Section~\ref{sec:relax}).  Optimal
clusters have diameter at most two in $G^+$ (Lemma~\ref{lem:half}).  We show
that the metric LP restricted to pairs at distance at most two, with all pairs
at distance at least three fixed to ``separated'' and the triangle
inequalities through them kept, is still a relaxation.  Moreover, every
feasible point extends to a feasible point of the full metric LP with the same
cost (Theorem~\ref{thm:sparse}).  Hence the $2.5$- and $2.06$-approximate
LP-rounding algorithms~\citep{AilonCN08,ChawlaMSY15} keep their guarantees,
and the relaxation has $O(\sum_v\deg(v)^2)$ variables instead of
$\binom{n}{2}$.  We strengthen it with star inequalities and \emph{local
subgraph inequalities} (Proposition~\ref{prop:subgraph}).  On real graphs
these raise the bound substantially.  On \texttt{ca-GrQc}, the gap between
the metric LP and the best known solution is $23\%$; with these inequalities
it falls below $0.2\%$.

\item \textbf{Anytime dual bounds} (Section~\ref{sec:dual}).  We compute a
lower bound by dual block-coordinate ascent.  Each block is a vertex set whose
LP is solved exactly by cutting planes, with costs equal to the current
reduced costs.  The bound is valid after every step, never decreases, and
stores only the reduced cost of each pair and the rows with positive
multipliers (Theorem~\ref{thm:anytime}).  Any packing of induced stars or bad
triangles is a feasible starting point.  We compute one on the sparse support with a
ruin-and-recreate local search, using $O(|P|)$ instead of $\Theta(n^2)$
memory.  On \texttt{ca-HepTh}, after 46\,s it is within $0.02\%$ of the
root bound of the KaPoCE branch-and-bound, which needs 277\,s; together with
the block LPs, our bound exceeds it.  Blocks come from randomised
graph partitions and from regions where the current optimality gap is large.

\item \textbf{Gap decomposition and local certificates}
(Section~\ref{sec:gap}).  For any clustering, the difference between its cost
and the dual bound splits exactly into non-negative terms attached to pairs
and constraints (Theorem~\ref{thm:gap}).  As a consequence, re-optimising a
vertex set $B$ can improve the cost by at most the gap mass touching $B$
(Corollary~\ref{cor:local}).  This gives a principled way to choose
neighbourhoods for a large-neighbourhood search with exact sub-MIPs, and to
skip neighbourhoods that are certified optimal.

\item \textbf{CertiFlip} (Section~\ref{sec:algo}).  We implement the iterated
flipping local search of \citet{CohenAddadLPTYZ24}, with cluster-insertion
moves generated by neighbourhood cleaning and multilevel refinement.  Together
with the three components above this yields an algorithm whose expected cost
is at most $3\,\OPT$, and at most $2.06\,\OPT$ when a single-block relaxation
is solved to optimality.  It returns a certified lower bound with every
solution (Theorem~\ref{thm:certiflip}).

\item \textbf{Exact recombination and PXMem} (Section~\ref{sec:pxmem}).  The
blocks of the overlap of two clusterings are independent: taking the better
parent in each block gives a clustering that is never worse than either
parent (Theorem~\ref{thm:px}).  Together with the contraction of critical
cliques (Lemma~\ref{lem:twins}) this yields PXMem, a memetic solver that keeps
the guarantee of $3$ in expectation (Theorem~\ref{thm:pxmem}) and beats
KaPoCE with statistical significance on four of eleven SNAP graphs.

\item \textbf{Machine-checked proofs} (Section~\ref{sec:lean}).  The
combinatorial lemmas and theorems of the paper, the incremental cost formulas
of the implementation and the certificate and approximation arguments are
proved in Lean~4 with Mathlib.

\item \textbf{Evaluation pipeline and experiments} (Section~\ref{sec:exp}).
We release an open-source pipeline with ten solvers and five lower bounds.
On the 200 PACE~2021 exact-track instances, whose optima are mostly known, the
CertiFlip bound proves optimality on 111 of the 173 solved instances.  The
root bounds of the state-of-the-art branch-and-bound prove it on 79.  The
CertiFlip bound also improves the published lower bounds of two open
instances.  We also use fifteen SNAP graphs with up to $1.1\cdot 10^6$ vertices.  On the
eleven whose distance-two support fits into memory (up to $3\cdot10^5$
vertices and $10^6$ edges), the CertiFlip bound certifies the best known
solutions within $0.2\%$ to $9.2\%$ of optimal.  Previously scalable bounds
certified only factors of $1.17$ to $1.81$ on the same graphs.
\end{enumerate}

\section{Preliminaries and related work}\label{sec:prelim}

Let $G^+=(V,\Ep)$ with $n=|V|$ and $m=|\Ep|$; every pair not in $\Ep$ is
negative.  For a clustering $\mathcal{C}$ let $x^{\mathcal C}_{uv}=1$ if $u,v$
are in different clusters and $0$ otherwise.  Then
\[
\cost(\mathcal C)=\sum_{uv\in\Ep}x^{\mathcal C}_{uv}+\sum_{uv\notin\Ep}(1-x^{\mathcal C}_{uv}).
\]
$N(v)$ denotes the positive neighbourhood and $d(u,v)$ the distance in $G^+$.
The \emph{metric LP}~\citep{CharikarGW05,AilonCN08} relaxes $x$ to $[0,1]$
subject to the triangle inequalities $x_{uv}\le x_{uw}+x_{wv}$ for all
triples.  A \emph{bad triangle} is a triple $u,w,v$ with $uw,wv\in\Ep$ and
$uv\notin\Ep$; every clustering makes at least one mistake on it.

\paragraph{Approximation algorithms.}  Pivot repeatedly picks a uniformly
random unclustered vertex and clusters it with its unclustered positive
neighbours.  Its expected cost is at most three times the value of the
fractional bad-triangle packing LP~\citep{AilonCN08}.  LP-based pivoting joins
$v$ to the pivot $u$ with probability $1-f(x_{uv})$.  The choice $f(x)=x$
gives $2.5$~\citep{AilonCN08}.  \citet{ChawlaMSY15} use $f^-(x)=x$ on negative
pairs and, on positive pairs, $f^+(x)=0$ for $x<a$,
$((x-a)/(b-a))^2$ for $a\le x<b$ and $1$ for $x\ge b$, with
$a=0.19$, $b=0.5095$.  Both guarantees hold against the LP value of \emph{any}
feasible metric $x$, a fact we use below.  For the history of ratios,
parallel, streaming and dynamic algorithms and weighted variants
see~\citep{CohenAddadLPTYZ24,CaoCLLLNTVYZ25,GarciaSorianoS26,
cohenaddad2021constant,behnezhad2022almost3,DemaineEFI06}.  Relevant here is
the local search of \citet{CohenAddadLPTYZ24}.  Its only move is a
\emph{cluster insertion} $\mathcal C+S$: the vertices of $S$ leave their
clusters and form a new cluster.  A clustering that no insertion improves is a
$2$-approximation.  Increasing the weight of the positive pairs cut by the
current solution (a \emph{flip}), re-optimising, and combining the last three
solutions with a $3$-way pivot yields $2-2/13$.

\paragraph{Exact and heuristic solvers.}  Cluster editing admits linear
kernels~\citep{Guo09kernel,CaoC12kernel} and practical branch-and-bound
algorithms~\citep{BockerBK11exact,BlasiusEtAl22sea}.  The latter use packings
of stars and induced paths as lower bounds, and mark pairs at distance at
least three as forbidden.  \citet{BastosOPSMP16} note that such pairs are
separated in every optimal solution.  Cutting-plane methods for the clique
partitioning polytope go back to \citet{GrotschelW89}, with recent work on
facets and redundancy~\citep{irmai2024chorded,koshimura2022concise}.  For
sparse weighted graphs (multicut), dual decomposition and message passing
give fast lower bounds~\citep{SwobodaA17multicut,AbbasS22rama,lange2018partial}.
On complete graphs, however, the number of variables is quadratic.  Our
distance-two support removes exactly this obstacle.  Local search heuristics
include vertex moves~\citep{elsner2009bounding}, Louvain and Leiden with the
constant Potts model at resolution $1/2$, which optimises the same
objective~\citep{veldt2018lambdacc,BlondelGLL08,TraagWvE19}, and memetic
multilevel schemes~\citep{hausberger2025scc,BlasiusEtAl21kapoceHeur}.

\section{Distance-two relaxations}\label{sec:relax}

\begin{lemma}\label{lem:half}
Let $\mathcal C$ be an optimal clustering, $S\in\mathcal C$ and $v\in S$.  Then
$|N(v)\cap S|\ge (|S|-1)/2$.  Consequently any two vertices of $S$ are
adjacent or have a common neighbour in $S$.  In particular every pair at
distance at least three is separated.
\end{lemma}
\begin{proof}
Moving $v$ to a singleton changes the cost by $|N(v)\cap S|-(|S|-1-|N(v)\cap
S|)$, which must be non-negative.  For non-adjacent $u,v\in S$ the sets
$N(u)\cap S$ and $N(v)\cap S$ lie in $S\setminus\{u,v\}$, which has $|S|-2$
elements, and have total size at least $|S|-1$, so they intersect.
\end{proof}

Let $\Nt=\{uv\notin\Ep: N(u)\cap N(v)\ne\emptyset\}$, $P=\Ep\cup\Nt$ and
$F=\binom V2\setminus P$ (the \emph{far} pairs).  Note
$|P|\le m+\sum_{w}\binom{\deg(w)}{2}$.

\begin{definition}[distance-two relaxation]
$\mathrm{LP}_P$ has a variable $x_e\in[0,1]$ for every $e\in P$ and minimises
$\sum_{e\in\Ep}x_e+\sum_{e\in\Nt}(1-x_e)$ subject to
\begin{align}
x_{uv}&\le x_{uw}+x_{wv} &&\text{for all } u,v,w \text{ with } uv,uw,wv\in P,\label{eq:t1}\\
x_{uw}+x_{wv}&\ge 1 &&\text{for all } u,v,w \text{ with } uw,wv\in P,\ uv\in F.\label{eq:t2}
\end{align}
\end{definition}

\begin{theorem}\label{thm:sparse}
Let $\mathrm{LP}_{\mathrm{std}}$ denote the value of the metric LP.
\begin{enumerate}
\item[(a)] $\mathrm{LP}_{\mathrm{std}}\le \mathrm{LP}_P\le\OPT$.
\item[(b)] If $x$ is feasible for $\mathrm{LP}_P$, then $\bar x$, defined as $x$ on
$P$ and $1$ on $F$, is feasible for the metric LP and has the same cost.
\item[(c)] Let $\mathcal A$ be a randomised rounding algorithm with
$\mathbb E[\cost(\mathcal A(\bar x))]\le\alpha\,\mathrm{cost}(\bar x)$ for every
feasible metric $\bar x$.  Applied to an optimal solution of $\mathrm{LP}_P$ it
is an $\alpha$-approximation.  This holds in particular for $\alpha=2.5$
(ACN) and $\alpha=2.06$ (CMSY).  With pivot-based rounding, only pairs in
$P$ can be joined, so the rounding takes $O(n+|P|)$ time.
\end{enumerate}
\end{theorem}
\begin{proof}
(b) Consider a triple $u,v,w$ and the number of far pairs among $uv,uw,wv$.
With none, all three triangle inequalities are rows~\eqref{eq:t1}.  With
exactly one, say $uv$, the inequality $\bar x_{uv}=1\le x_{uw}+x_{wv}$ is
row~\eqref{eq:t2}.  The other two inequalities, such as
$x_{uw}\le 1+x_{wv}$, hold because $x\le 1$.  With two or three far pairs, each
inequality has a right-hand side of at least $1$.  Far pairs are negative and
$\bar x=1$ on them, so they contribute $0$ to the cost.  (a) The first
inequality follows from (b).  For the second, by Lemma~\ref{lem:half} an
optimal clustering separates all far pairs.  Its cut vector restricted to $P$
therefore satisfies~\eqref{eq:t1} and~\eqref{eq:t2}, because it is a metric
with $x=1$ on $F$, and it has cost $\OPT$.  (c) follows from (a) and (b): the
expected cost is at most $\alpha\,\mathrm{cost}(\bar x)=\alpha\,\mathrm{LP}_P
\le\alpha\,\OPT$.  Pivot-based rounding joins $v$ to the pivot $u$ with
probability $1-f(x_{uv})$, and $f(1)=1$ for both rounding functions.
\end{proof}

The rows~\eqref{eq:t1}--\eqref{eq:t2} are generated lazily.  A violated row
needs $x_{uw}+x_{wv}<1$, so separation around a vertex $w$ only has to look at
pairs of ``close'' neighbours in $P$, which keeps it fast in practice.

\paragraph{Stronger inequalities.}  The metric LP has integrality gap $2$,
already on a star.  We add two families that are valid for every clustering
separating all far pairs, and hence for every optimal clustering.

\begin{proposition}\label{prop:subgraph}
Let $\mathcal C$ be a clustering in which no far pair shares a cluster.
\begin{enumerate}
\item[(i)] (star inequalities; the case $|S|=1$ of the $2$-partition
inequalities of \citet{GrotschelW89}) For every $v$ and $T\subseteq V\setminus\{v\}$, with
$y=1-x^{\mathcal C}$,
\[\textstyle\sum_{t\in T}y_{vt}-\sum_{\{t,t'\}\subseteq T,\ tt'\in P}y_{tt'}\le 1.\]
\item[(ii)] (local subgraph inequalities) For every $H\subseteq V$,
\[\textstyle\sum_{e\in\Ep(H)}x^{\mathcal C}_e+\sum_{e\in\Nt(H)}(1-x^{\mathcal C}_e)\ \ge\ \OPT(H),\]
where $\OPT(H)$ is the optimum of the instance induced by $H$ and $\Ep(H)$,
$\Nt(H)$ are the pairs of $\Ep$ and $\Nt$ inside $H$.
\end{enumerate}
\end{proposition}
\begin{proof}
(i) If the cluster of $v$ contains $k$ vertices of $T$, the first sum is $k$
and the second is at least $\binom k2$.  Pairs inside $T$ that are not in $P$
are far, so they are separated and contribute $0$ to the second sum.
Finally $k-\binom k2\le1$.  (ii) Restricted to $H$, $\mathcal C$ is a clustering
of $H$.  Its mistakes inside $H$ number at least $\OPT(H)$.  The far pairs
inside $H$ are separated and cost nothing, so all its mistakes are counted
on the left-hand side.
\end{proof}

Bad triangles ($\OPT=1$) and induced stars with $k$ independent leaves
($\OPT=k-1$) are local subgraph inequalities.  So are chordless cycles $C_k$,
for which $\OPT(C_k)=\lceil k/2\rceil$ while the metric LP only certifies
$k/2$ (put $x=1/2$ on the cycle edges).  We separate stars greedily around
each vertex.  For (ii) we grow $H$ ($|H|\le 8$) by breadth-first search along
fractional pairs and compute $\OPT(H)$ by enumerating the at most $4140$
set partitions of $H$.

\section{Anytime dual bounds}\label{sec:dual}

Write the relaxation as $\min\{K+c^\top x: Ax\ge b,\ 0\le x\le 1\}$ over
$x\in\R^P$, with $c_e=1$ on $\Ep$, $c_e=-1$ on $\Nt$ and $K=|\Nt|$.  The rows
of $(A,b)$ are any valid inequalities from Section~\ref{sec:relax}.

\begin{theorem}\label{thm:anytime}
For every $y\ge 0$ let $r=c-A^\top y$ and
\begin{equation}\label{eq:lb}
\LB(y)=K+b^\top y+\sum_{e\in P}\min\{0,r_e\}.
\end{equation}
\begin{enumerate}
\item[(a)] $\LB(y)\le\OPT$.
\item[(b)] Let $B\subseteq V$, let $R_B$ be a set of rows supported on pairs
with both endpoints in $B$, and let $\bar r$ be the reduced costs after
removing the contribution of the current multipliers on $R_B$.  Replace these
multipliers by an optimal dual solution of
$\min\{\bar r^\top x_B: A_Bx_B\ge b_B,\ 0\le x_B\le 1\}$.  This does not
decrease $\LB$.
\item[(c)] If $B=V$ and $R_B$ contains all rows of $\mathrm{LP}_P$, one block
step gives $\LB=\mathrm{LP}_P$.
\end{enumerate}
\end{theorem}
\begin{proof}
(a) Let $\mathcal C$ be optimal.  By Lemma~\ref{lem:half} it separates far
pairs, so $\OPT=K+c^\top x^{\mathcal C}$ with $x^{\mathcal C}\in\{0,1\}^P$
satisfying $Ax^{\mathcal C}\ge b$.  Then
$\OPT=K+r^\top x^{\mathcal C}+y^\top Ax^{\mathcal C}\ge K+\sum_e\min\{0,r_e\}+y^\top b$.
(b) By LP duality the block LP equals
$\max_{y_B\ge0}\ b_B^\top y_B+\sum_{e\in P_B}\min\{0,(\bar r-A_B^\top
y_B)_e\}$.  The terms of~\eqref{eq:lb} outside the block do not change, and
the old multipliers are feasible for this maximisation.  (c) is LP duality
for $\mathrm{LP}_P$ with the box constraints.
\end{proof}

\paragraph{Warm start by star packings.}  Every pair-disjoint family of
induced stars (a centre $v$ with $k\ge2$ pairwise non-adjacent positive
neighbours) is a feasible dual solution: put $y=1$ on the corresponding
local subgraph inequalities of Proposition~\ref{prop:subgraph}(ii), which in
$x$-form read $\sum_t x_{vt}-\sum_{tt'}x_{tt'}\ge (k-1)-\binom k2$.  Then
$\LB(y)$ equals the packing value $\sum(k-1)$.  Indeed, a centre pair has
$c_e=1$ and coefficient $1$, and a leaf pair has $c_e=-1$ and coefficient
$-1$, so every packed pair has reduced cost $0$.  Unpacked pairs keep
$r_e=c_e$.  Hence $\LB=K+\sum_S\big((k_S-1)-\binom{k_S}2\big)-|\Nt\setminus
\text{packed}|$, and $K-|\Nt\setminus\text{packed}|=\sum_S\binom{k_S}{2}$.
Fractional packings of bad triangles, as computed by the Garg--K\"onemann
scheme~\citep{GargK07,Fleischer00}, are installed in the same way on the
rows $x_{uw}+x_{wv}-x_{uv}\ge0$.
Packings are the bounds used inside the KaPoCE branch-and-bound
algorithm~\citep{BlasiusEtAl22sea}.  Its implementation stores the instance
as a dense matrix, so it applies only to graphs with at most about $10^4$
vertices.  We build a greedy packing on the sparse support, processing
centres by decreasing degree, and improve it by \emph{ruin and recreate}.
(On dense instances large stars waste pairs, since a star with $k$ leaves
uses $k+\binom k2$ pairs for a value of $k-1$.  There we also try a packing
that first packs bad triangles maximally and then extends stars on free
pairs, as well as the fractional triangle packing, and install the best of
the candidates.)
Each step picks a random vertex, removes all stars centred in a random
connected region of about 12 vertices around it, and repacks the freed
pairs greedily in a random order.  A step is kept iff the value does not
decrease.  The packing is installed as the initial multipliers.  Subsequent
block steps release the packed rows that lie inside a block and re-optimise
them fractionally together with newly separated rows.  By
Theorem~\ref{thm:anytime}(b), the final bound is never below the packing
value.

The bound~\eqref{eq:lb} needs only the vector $r$ and the scalar $b^\top y$.
We store the rows with positive multipliers so that a later block containing
them can re-optimise them.  Blocks of the same partition are independent.
In a \emph{sweep} we compute a partition into parts of about $2500$ vertices
with METIS~\citep{KarypisK98} on a randomly relabelled copy of $G^+$, so that
successive sweeps cut different pairs.  Each part is solved with the dual
simplex method of HiGHS~\citep{HuangfuH18}, with lazily separated triangle,
star and local subgraph rows.  After the sweeps stall we switch to
\emph{gap-guided blocks}, grown around the vertices with the largest share of
the gap defined next.  Because~\eqref{eq:lb} is valid for any $y\ge0$, the
bound remains correct even when a block LP hits its time limit.  We only
clamp the dual values to be non-negative.

\section{Gap decomposition and local certificates}\label{sec:gap}

\begin{theorem}\label{thm:gap}
Let $y\ge0$, $r=c-A^\top y$, and let $\mathcal C$ be a clustering that separates all
far pairs.  Then
\begin{equation}\label{eq:gap}
\cost(\mathcal C)-\LB(y)=\sum_{e\in P}\underbrace{\big(r_ex^{\mathcal C}_e-\min\{0,r_e\}\big)}_{g_e\ge0}
+\sum_i\underbrace{y_i\,(a_i^\top x^{\mathcal C}-b_i)}_{s_i\ge0}.
\end{equation}
\end{theorem}
\begin{proof}
$\cost(\mathcal C)=K+c^\top x^{\mathcal C}=K+r^\top x^{\mathcal C}+y^\top
Ax^{\mathcal C}$.  Subtract~\eqref{eq:lb}.  Each $g_e\ge0$ because
$x_e\in\{0,1\}$, and each $s_i\ge0$ by the validity of row $i$.
\end{proof}

A pair term vanishes iff the clustering agrees with the sign of the reduced
cost.  A row term vanishes iff the row is tight or has multiplier $0$.
Assigning each term to the endpoints of its pairs gives a \emph{gap map} over
the vertices.

\begin{corollary}[local certificate]\label{cor:local}
For $B\subseteq V$ let $\Gamma_B(\mathcal C)$ be the sum of the terms
of~\eqref{eq:gap} whose pair, or one of whose row pairs, has an endpoint in
$B$.  Let $\mathcal C'$ be any clustering that separates far pairs and
coincides with $\mathcal C$ on all pairs with both endpoints outside $B$.  Then
$\cost(\mathcal C)-\cost(\mathcal C')\le\Gamma_B(\mathcal C)$.  In particular,
if $\Gamma_B(\mathcal C)<1$ then no re-clustering of $B$ improves $\mathcal C$.
\end{corollary}
\begin{proof}
Apply~\eqref{eq:gap} to both clusterings.  Terms not touching $B$ coincide,
and the terms of $\mathcal C'$ are non-negative.  Costs are integers.
\end{proof}

\paragraph{Gap-guided LNS.}  Large-neighbourhood search~\citep{Shaw98} fixes
the clustering outside a vertex set $B$ ($|B|\le 60$) and re-optimises $B$
exactly.  Let $\mathcal K$ be the outside clusters with a positive edge into
$B$.  Re-clusterings of $B$ correspond to clusterings of $B$ together with one
super node per $K\in\mathcal K$, where super nodes must stay apart.  The cost is
$\sum_{uv}[uv\in\Ep]x_{uv}+[uv\notin\Ep](1-x_{uv})$ for $u,v\in B$, plus
$w_{vK}x_{vK}+(|K|-w_{vK})(1-x_{vK})$ for each $v\in B$, where $w_{vK}$ counts
the positive edges between $v$ and $K$.  We solve this MIP with HiGHS and add
triangle inequalities lazily until the solution is transitive.  To keep
Corollary~\ref{cor:local} applicable, $v$ may join $K$ only if it is within
distance two of every member of $K$, and far pairs inside $B$ are fixed apart.
Neighbourhoods are grown around high-gap vertices, either by breadth-first
search or as a union of whole clusters.  Neighbourhoods with $\Gamma_B<1$ are
skipped because they are certified optimal.

\section{The CertiFlip algorithm}\label{sec:algo}

\begin{algorithm}[t]
\caption{CertiFlip$(G^+,T)$}\label{alg:certiflip}
\begin{algorithmic}[1]
\State $\mathcal C\gets\textsc{Pivot}(G^+)$ with a uniformly random order
\State $\mathcal C\gets\textsc{IteratedFlip}(\mathcal C)$ \Comment{insertion local search with flips}
\State build the distance-two support $P$
\State $y\gets$ star packing with ruin-and-recreate local search \Comment{feasible dual}
\State $y\gets\textsc{BlockAscent}(P,y,\mathcal C)$ \Comment{METIS sweeps, then gap-guided blocks}
\If{one block covered $V$ and separation converged}
  \State $\mathcal C\gets$ best of $\mathcal C$ and \textsc{IteratedFlip}(\textsc{CMSY-Round}$(x^\star)$) over several roundings
\EndIf
\State $\mathcal C\gets\textsc{GapLNS}(\mathcal C,y)$ until time $T$
\State \Return $\mathcal C$ and $\lceil\LB(y)\rceil$ \Comment{lines 4--5 share half of the budget $T$}
\end{algorithmic}
\end{algorithm}

\paragraph{Insertion local search.}  For every root $r$ we build one candidate
insertion $S(r)$ by \emph{cleaning} $N[r]$.  Starting from $S=N[r]$, we
repeat a few sweeps over $v\in N[r]\setminus\{r\}$: we put $v$ in $S$ iff
joining $S$ costs $v$ less than staying in the remainder of its current
cluster.  Both costs are computed exactly from counts of weighted positive
neighbours.  We then compute the exact change of the objective in
$O(\sum_{v\in S}\deg v)$ time and apply the insertion if it is negative.
Rounds of insertions alternate with a multilevel (Louvain-type) vertex-move
local search.  That search works on the weighted objective, moves super nodes
at coarse levels and hence also merges clusters.  \textsc{IteratedFlip}
follows \citet{CohenAddadLPTYZ24}.  Positive weights take values in
$\{1,1.5,2\}$ (we double all weights to stay in integers).  Each round flips
the pairs cut by the previous solution, re-optimises, flips and re-optimises
again, and combines the last three solutions with their $3$-way pivot.  The
best candidate under the original objective is returned.

\begin{theorem}\label{thm:certiflip}
Let $\mathcal C$ and $\LB$ be the output of Algorithm~\ref{alg:certiflip}.
\begin{enumerate}
\item[(a)] $\mathbb E[\cost(\mathcal C)]\le3\,\OPT$.
\item[(b)] $\LB\le\OPT$, so $\cost(\mathcal C)/\lceil\LB\rceil$ is an upper
bound on the approximation ratio achieved on the instance.
\item[(c)] If line~7 is executed, $\mathbb E[\cost(\mathcal C)]\le 2.06\,\OPT$.
\end{enumerate}
\end{theorem}
\begin{proof}
Every step after line~1 returns a clustering no worse than its input: the
local searches only apply improving moves, \textsc{IteratedFlip} returns the
best of its candidates, which include the local optimum reached from its input,
and \textsc{GapLNS} accepts only improvements.  (a) then follows from the
Pivot guarantee~\citep{AilonCN08}.  (b) is Theorem~\ref{thm:anytime}(a) and
the integrality of $\OPT$.  (c) If one block covered $V$ and separation
converged, the block LP primal $x^\star$ satisfies all
rows~\eqref{eq:t1}--\eqref{eq:t2}.  Its cost equals the block LP value, which
is at most $\mathrm{LP}_P$ plus the value added by the extra valid rows, and
hence at most $\OPT$.  By Theorem~\ref{thm:sparse}(b,c),
$\mathbb E[\cost(\textsc{CMSY-Round}(x^\star))]\le 2.06\,\OPT$.  The remaining
steps are monotone.
\end{proof}

\begin{remark}
If the insertion local search reached an optimum with respect to
\emph{all} cluster insertions, \citet{CohenAddadLPTYZ24} would give a factor of
$2$, and $2-2/13$ with enough flipping rounds.  Our candidate generation
explores a polynomial subfamily of insertions, so we do not claim these
factors.  The certificate of Theorem~\ref{thm:certiflip}(b), however, is
valid in every case.
\end{remark}

\paragraph{Running time.}  Pivot takes $O(n+m)$ time.  One round of insertions
takes $O(\sum_v\sum_{u\in N(v)}\deg u)=O(\sum_u\deg(u)^2)$, the same order as
the size of $P$ and the number of bad triangles.  The support takes
$O(\sum_w\deg(w)^2)$ time and memory.  The block LPs dominate the running
time, and the time spent on them is set by the budget $T$.

\section{Exact recombination: the PXMem solver}\label{sec:pxmem}
\input{pxmem}

\section{Machine-checked proofs}\label{sec:lean}
The correctness of the solvers rests on a number of short combinatorial
arguments and on exact formulas that the implementation evaluates millions of
times per second.  An error in a sign or in a boundary case would not crash
the code; it would silently produce worse or invalid results.  We therefore
formalised part of the mathematics in Lean~4~\citep{deMouraU21lean4} with
Mathlib~\citep{mathlib2020}.  The development has about 1{,}800 lines in
fourteen files and is part of the repository (\texttt{lean/CCProofs}).  Every
theorem named in Table~\ref{tab:lean} is checked by the Lean kernel without
\texttt{sorry} and depends only on the standard axioms \texttt{propext},
\texttt{Classical.choice} and \texttt{Quot.sound}.

The formal model is the one of Section~\ref{sec:prelim}: a finite vertex set,
a simple graph $G^+$ and a labelling $c:V\to\alpha$ into an arbitrary type
with decidable equality.  The cost is summed over ordered pairs, so every
count appears multiplied by two (\texttt{cost\_eq\_two\_mul\_card} relates it
to the unordered count).  The weighted instance has integer node sizes and
symmetric weights, as in~\eqref{eq:wcost}.

\begin{table}[!tbp]
\centering
\caption{Coverage of the formal development.  ``Full'': the statement of the
paper is proved as stated (up to the ordered-pair convention); ``partial'':
the part named is proved; ``reduction'': proved from the cited guarantee
taken as a hypothesis; ``--'': not formalised.}\label{tab:lean}
\small
\begin{tabular}{p{0.33\linewidth}p{0.12\linewidth}p{0.47\linewidth}}
\toprule
result & coverage & Lean theorems \\
\midrule
Lemma~\ref{lem:half} & partial & far pairs are separated by every optimal clustering: \texttt{optimal\_separates\_far}, \texttt{sepFar\_of\_optimal} \\
Theorem~\ref{thm:sparse} & -- & \\
Proposition~\ref{prop:subgraph}(i) & full & \texttt{star\_row\_valid} (star rows in $x$-form, far-separating clusterings) \\
Proposition~\ref{prop:subgraph}(ii) & partial & bad triangles and induced stars: \texttt{bad\_triangle\_costIn}, \texttt{star\_bound}; packings: \texttt{packing\_bound} \\
rows \eqref{eq:t1}--\eqref{eq:t2} & full & \texttt{sep\_triangle}, \texttt{sep\_far\_triangle} \\
Theorem~\ref{thm:anytime}(a) & full & \texttt{far\_dual\_bound}, \texttt{far\_dual\_le\_opt} (bound on $P$ for rows valid for far-separating clusterings, $\le$ the cost of every clustering); \texttt{cost\_eq\_supp} \\
Theorem~\ref{thm:anytime}(b,c), Theorem~\ref{thm:gap}, Corollary~\ref{cor:local} & -- & \\
Theorem~\ref{thm:certiflip} & (a) reduction, (b) full, (c) -- & (a) as Theorem~\ref{thm:pxmem}; (b) is Theorem~\ref{thm:anytime}(a) \\
Lemma~\ref{lem:twins} & full & \texttt{twins\_split\_improvable}, \texttt{optimal\_keeps\_twins} \\
Pivot keeps twins together & partial & one round: \texttt{pivot\_round\_twins} \\
contraction~\eqref{eq:wcost} & partial & cost identity: \texttt{cost\_contract} \\
move and swap formulas~\eqref{eq:move} & full & \texttt{cost\_move}, \texttt{costW\_move}, \texttt{costW\_swap}, \texttt{cost\_update} \\
Theorem~\ref{thm:px} & full & any pairwise-additive cost: \texttt{crossoverP\_le}; unweighted: \texttt{crossover\_le}; weighted: \texttt{crossoverW\_le}; relabelling: \texttt{cost\_comp\_injective}, \texttt{costW\_comp\_injective} \\
Theorem~\ref{thm:pxmem} & reduction & \texttt{reach\_le}, \texttt{output\_le}, \texttt{search\_guarantee}: non-worsening replacements keep a member no worse than the seed; the Pivot guarantee~\citep{AilonCN08} is a hypothesis \\
\bottomrule
\end{tabular}
\end{table}

\paragraph{Scope.}  The formalisation covers the mathematics, not the Python
and Numba code, and not all results of the paper (Table~\ref{tab:lean}).
Three points connect it to the numbers we report.  First, the lower bounds
are computed in floating point, but they are certified independently: the
solver writes out every row with a positive multiplier, and a separate checker
(\texttt{experiments/check\_certificate.py}), which shares no code with the
solver, (i)~checks that every pair of every row lies in $P$, (ii)~checks the
validity of every row for all far-separating clusterings, by enumerating the
far-separating partitions of its vertex set for rows on at most nine vertices
and by the closed form of \texttt{star\_row\_valid} for star rows, (iii)~rounds
the multipliers down to multiples of $2^{-30}$, which keeps them non-negative,
and (iv)~evaluates $b^\top y+\sum_{p}\bigl(\min\{0,c_p-(A^\top y)_p\}-\min\{0,c_p\}\bigr)$
in exact integer arithmetic.  This is the left-hand side of
\texttt{far\_dual\_bound}, so every bound we report as certified is a proved
lower bound on $\OPT$ up to the correctness of this short checker, not of the
solver.  Second, \texttt{star\_row\_valid} and \texttt{sep\_far\_triangle} hold
only for clusterings that separate far pairs; they enter the bound through
\texttt{far\_dual\_le\_opt}, which uses that optimal clusterings do so.  Third,
the code is tested against the formal statements: exact costs of random
clusterings on contracted instances, the tracked cost against a recomputation
at the end of annealing runs with many moves, swaps and rollbacks, and brute
force on small graphs.  These tests check final states, not every step.


\section{Experiments}\label{sec:exp}
\subsection{Setup}

\paragraph{Implementation.}  All our algorithms are implemented in Python
with Numba-compiled kernels in the package \texttt{ccbench}.  LPs and MIPs are
solved with HiGHS~1.15~\citep{HuangfuH18} and graph partitions are computed
with METIS~\citep{KarypisK98}.  Experiments ran on a 4-core virtual machine
with 15\,GB of memory.  Each run used one thread.  Four runs executed in parallel on the PACE
instances and three on the SNAP graphs.

\paragraph{Algorithms.}  We compare the following.
\emph{Pivot}~\citep{AilonCN08}.
\emph{Pivot-50}: the best of 50 Pivot runs.
\emph{MatchFlipPivot}~\citep{veldt2022stc}.
\emph{Vote}: best-improvement vertex moves from singletons, in the spirit
of~\citep{elsner2009bounding}.
\emph{Louvain-CC}: our multilevel implementation of the vertex-move local
search.
\emph{Leiden-CPM}: Leiden with the constant Potts model at resolution
$1/2$~\citep{TraagWvE19}, which optimises the correlation clustering
objective~\citep{veldt2018lambdacc}.
\emph{Insertion LS}: our cluster-insertion local search.
\emph{IteratedFlip}: five flipping rounds, following~\citep{CohenAddadLPTYZ24}.
\emph{KaPoCE}: the PACE~2021 winner of both tracks, heuristic
version~\citep{BlasiusEtAl21kapoceHeur}.
\emph{CertiFlip}: Algorithm~\ref{alg:certiflip}.
\emph{PXMem}: Algorithm~\ref{alg:pxmem} (Section~\ref{sec:h2h}).
KaPoCE and CertiFlip receive the same wall-clock budget $T$.  The other
methods run to completion, with three seeds each.

\paragraph{Lower bounds.}  We report the following bounds.
\emph{Greedy packing}: a maximal pair-disjoint packing of bad triangles, i.e.\
the matching bound of MatchFlipPivot.
\emph{Triangle packing}: a fractional packing by the Garg--K\"onemann
scheme~\citep{GargK07,Fleischer00} with $\varepsilon=0.05$, normalised by its
maximum load, hence always valid.  Its value approximates the covering LP
over bad triangles.
The \emph{CertiFlip bound}: Theorem~\ref{thm:anytime}.
On PACE instances we also use the root \emph{star-packing} and
\emph{$P_3$-packing} bounds of the KaPoCE branch-and-bound, as published by
its authors~\citep{BlasiusEtAl22sea}.  All reported bounds are rounded up to
integers.

\paragraph{Instances.}  (i)~The 200 instances of the PACE~2021 exact
track~\citep{KellerhalsKNZ21pace} ($n\le 620$).  173 of them have a known
optimum and 27 are open.  (ii)~Fifteen SNAP graphs, used as complete
instances with $\Ep$ given by the edges (for signed networks, by the positive
edges): collaboration networks \texttt{ca-*}, \texttt{email-Enron},
\texttt{loc-Brightkite}, \texttt{soc-Epinions}, the positive parts of
Slashdot, Epinions and the two Bitcoin trust networks~\citep{leskovec2010signedchi,kumar2016bitcoin},
and the community graphs \texttt{com-Amazon}, \texttt{com-DBLP} and
\texttt{com-Youtube}~\citep{yang2012groundtruth}.  (iii)~The 200 instances of the PACE~2021 heuristic
track~\citep{KellerhalsKNZ21pace} (up to $2.0\cdot10^6$ vertices), including the
100 instances that were hidden during the challenge; they were used only for
the final comparison with KaPoCE.

\subsection{PACE 2021 exact track}

Table~\ref{tab:pace-primal} compares the solutions with the best known value
of each instance.  That value is the published optimum or upper
bound~\citep{BlasiusEtAl22sea}, or the best value found in our runs if lower.
KaPoCE and CertiFlip had a budget of $60$\,s.  The other methods ran to
completion, and we report the mean over three seeds.

\begin{table}[t]
\centering
\caption{PACE 2021 exact track (200 instances): primal quality.  ``Hit''
counts instances on which the best of the runs matches the best known value;
the gap is the relative excess of the mean cost over the best known value.}
\label{tab:pace-primal}
\begin{tabular}{lrrrr}
\toprule
Algorithm & best known hit & mean gap (\%) & max gap (\%) & time (s)\\
\midrule
Pivot & 5/42 & 57.906 & 227.08 & 0.06\\
Pivot-50 & 8/42 & 12.221 & 26.07 & 0.04\\
MatchFlipPivot & 1/42 & 143.090 & 362.50 & 0.07\\
Vote & 35/42 & 1.450 & 11.11 & 0.03\\
Louvain-CC & 37/42 & 0.563 & 5.56 & 0.04\\
Leiden-CPM & 40/42 & 0.332 & 3.17 & 0.37\\
Insertion LS & 39/42 & 0.193 & 3.02 & 0.05\\
IteratedFlip & 42/42 & 0.018 & 0.45 & 0.12\\
KaPoCE & 40/40 & 0.000 & 0.00 & 11.86\\
CertiFlip & 40/40 & 0.000 & 0.00 & 11.35\\
\bottomrule
\end{tabular}

\end{table}

\paragraph{Primal quality.}  Our implementation of the iterated flipping local
search matches the best known value on 187 of the 200 instances.  Its mean gap
is $0.018\%$ and a run takes three seconds on average.  Its approximation
guarantee is only $3$ in expectation, but in practice it is an order of
magnitude better than Louvain- or Leiden-type local search.  Those methods
optimise the same objective but reach only 153 and 158 instances, with mean
gaps of $0.54\%$ and $0.45\%$.  CertiFlip reaches 191 instances, and its
largest gap is $0.32\%$.  KaPoCE reaches all 200 instances.  This is not
surprising: the best known values come from its own branch-and-bound
runs~\citep{BlasiusEtAl22sea}, and it was engineered for this benchmark.
Pivot and MatchFlipPivot are far from optimal on these fairly dense
instances.  Among the 27 open instances, the best published upper bound of
\texttt{exact183} is 2789, and both KaPoCE and CertiFlip find solutions of
cost 2788 within the budget, so this improvement is shared by the two solvers.
We verified the CertiFlip solution independently, and it is included in the
repository.

\begin{table}[t]
\centering
\caption{PACE 2021 exact track: lower bounds on the 173 instances with known
optimum.  The last two rows are the root bounds of the KaPoCE
branch-and-bound as published by its authors.  ``LB $=$ OPT'' counts the
instances whose bound proves optimality.}
\label{tab:pace-bounds}
\begin{tabular}{lrrrr}
\toprule
Lower bound & instances & mean LB/OPT & min LB/OPT & LB $=$ OPT\\
\midrule
CertiFlip bound & 40 & 0.9934 & 0.951 & 31\\
Triangle packing (MWU) & 41 & 0.9502 & 0.814 & 13\\
Greedy packing & 41 & 0.8248 & 0.696 & 6\\
B\&B star packing & 173 & 0.9948 & 0.891 & 79\\
B\&B $P_3$ packing & 173 & 0.9742 & 0.813 & 59\\
\bottomrule
\end{tabular}

\end{table}

\paragraph{Lower bounds.}  Table~\ref{tab:pace-bounds} compares the bounds.
All CertiFlip bounds in this section are the values returned by the
independent checker of Section~\ref{sec:lean} on the exported
certificates; the checker accepted all 200 certificates.  The CertiFlip bound
(Theorem~\ref{thm:anytime}), computed within half of the
60\,s budget, proves optimality on 111 of the 173 instances.  The root
star-packing bound of the KaPoCE branch-and-bound proves optimality on 79,
its $P_3$-packing bound on 59, and the classical packings of bad triangles on
at most 30.  Our bound exceeds the star-packing bound on 60 instances, is
lower on 43 and equal on 70.  The comparison depends on density.  On the 35
instances with edge density below $0.1$, our bound attains $99.8\%$ of the
optimum on average, against $99.3\%$ for the star packing.  On denser
instances the star packing is better ($99.4\%$ to $99.7\%$ against
$98.0\%$ to $98.8\%$).  There the block LPs are too large to converge within the budget,
and the bound falls back to the best packing.

\begin{table}[t]
\centering
\caption{The 27 open PACE 2021 exact instances: published upper bound (UB)
and star-packing lower bound (LB)~\citep{BlasiusEtAl22sea}, our best
solution over all methods and runs, and the checked CertiFlip bound of one
60\,s run.  Bold marks improvements
over the published values.}
\label{tab:pace-open}
\small
\begin{tabular}{lrrrrrr}
\toprule
Instance & $n$ & $m$ & published UB & published LB & our UB & our LB\\
\midrule
exact042 & 100.0 & 1563.0 & 1055 & 1006 & 1055 & --\\
exact043 & 100.0 & 1624.0 & 984 & 968 & -- & --\\
exact051 & 120.0 & 2251.0 & 1619 & 1508 & -- & --\\
exact052 & 120.0 & 2448.0 & 1473 & 1431 & -- & --\\
exact060 & 124.0 & 2288.0 & 1491 & 1427 & -- & --\\
exact069 & 140.0 & 2915.0 & 2140 & 1994 & -- & --\\
exact070 & 140.0 & 3014.0 & 1792 & 1762 & -- & --\\
exact082 & 160.0 & 4015.0 & 2434 & 2397 & -- & --\\
exact083 & 160.0 & 4259.0 & 2939 & 2808 & -- & --\\
exact091 & 180.0 & 4602.0 & 3234 & 3072 & -- & --\\
exact092 & 180.0 & 6067.0 & 3563 & 3520 & -- & --\\
exact101 & 200.0 & 6731.0 & 4198 & 4135 & -- & --\\
exact102 & 200.0 & 6746.0 & 3925 & 3895 & -- & --\\
exact103 & 200.0 & 7081.0 & 4572 & 4381 & -- & --\\
exact105 & 200.0 & 12309.0 & 5320 & 5314 & -- & --\\
exact167 & 300.0 & 14191.0 & 9861 & 9458 & -- & --\\
exact169 & 300.0 & 15456.0 & 9561 & 9380 & -- & --\\
exact179 & 330.0 & 1370.0 & 633 & 620 & -- & --\\
exact180 & 330.0 & 2256.0 & 1081 & 1052 & -- & --\\
exact181 & 330.0 & 3253.0 & 1563 & 1513 & -- & --\\
exact182 & 330.0 & 4329.0 & 2159 & 2091 & -- & --\\
exact183 & 330.0 & 5613.0 & 2789 & 2717 & -- & --\\
exact184 & 330.0 & 7270.0 & 3684 & 3582 & -- & --\\
exact191 & 400.0 & 31201.0 & 18427 & 18288 & -- & --\\
exact195 & 620.0 & 3992.0 & 2529 & 2450 & -- & --\\
exact197 & 620.0 & 35188.0 & 15831 & 15782 & -- & --\\
exact198 & 620.0 & 91572.0 & 46764 & 46717 & -- & --\\
\bottomrule
\end{tabular}

\end{table}

Table~\ref{tab:pace-open} lists the open instances.  Our bound improves the
published lower bound on two of them: on \texttt{exact179} from 620 to 632,
so the optimum is now known to be 632 or 633, and on \texttt{exact180} from
1052 to 1068.  On
the other open instances, which are dense, the published star-packing bound
remains stronger.


\subsection{SNAP graphs}\label{sec:snap}

Table~\ref{tab:snap-primal} reports the best objective value per method on the
SNAP graphs.  KaPoCE and CertiFlip had $600$\,s.  CertiFlip spends half of its
remaining budget on the lower bound.  On the four graphs whose distance-two
support exceeds $3\cdot10^7$ pairs, the support does not fit into memory
together with three parallel runs.  There CertiFlip skips the bound and
spends the whole budget on the primal side.  MatchFlipPivot and the packing
bounds are not reported on these graphs.

\begin{table}[!tbp]
\centering
\caption{SNAP graphs: best objective value found (bold: best in the row).
All methods except KaPoCE and CertiFlip ran to completion; for these, the
best of three seeds is shown.  ``--'': support too large for memory
(MatchFlipPivot) or not finished within one hour (\texttt{com-Youtube}).}
\label{tab:snap-primal}
\scriptsize
\setlength{\tabcolsep}{3pt}
\resizebox{\textwidth}{!}{\begin{tabular}{lrrrrrrrrrrrr}
\toprule
Graph & $n$ & $m$ & Pivot & Pivot-50 & MatchFlipPivot & Vote & Louvain-CC & Leiden-CPM & Insertion LS & IteratedFlip & KaPoCE & CertiFlip\\
\midrule
\texttt{ca-GrQc} & 5242 & 14484 & 9811 & 8499 & 13068 & 6185 & 6158 & 6161 & 6080 & 6070 & \textbf{6050} & 6052\\
\texttt{ca-HepTh} & 9877 & 25973 & 22573 & 22573 & 29385 & 16255 & 16190 & 16193 & 15929 & 15877 & \textbf{15814} & 15857\\
\texttt{ca-HepPh} & 12008 & 118489 & 88732 & 63780 & 103445 & 46515 & 46477 & 46406 & 46222 & 46191 & \textbf{46133} & 46199\\
\texttt{ca-AstroPh} & 18772 & 198050 & 200307 & 184322 & 229913 & 126779 & 126633 & 126572 & 125746 & 125692 & \textbf{125559} & 125684\\
\texttt{ca-CondMat} & 23133 & 93439 & 87152 & 80134 & 110751 & 57856 & 57670 & 57611 & 57050 & 56998 & \textbf{56835} & 56980\\
\texttt{email-Enron} & 36692 & 183831 & 192960 & 185219 & 214677 & 151761 & 151611 & 151561 & 150906 & 150698 & \textbf{150430} & 150701\\
\texttt{loc-Brightkite} & 58228 & 214078 & 245496 & 236081 & 257088 & 178077 & 177773 & 177730 & 175734 & 175103 & \textbf{174695} & 175114\\
\texttt{soc-Epinions} & 75879 & 405740 & 496495 & 469475 & -- & 368339 & 368108 & 368094 & 366163 & 365580 & \textbf{365162} & 365601\\
\texttt{Slashdot+} & 75144 & 382882 & 469502 & 457004 & -- & 356033 & 355962 & 355915 & 353542 & 352349 & \textbf{351859} & 352367\\
\texttt{Epinions+} & 114467 & 592591 & 744964 & 669690 & -- & 525546 & 525237 & 525372 & 522976 & 522344 & \textbf{521825} & 522369\\
\texttt{BitcoinOTC+} & 5573 & 18591 & 20058 & 19456 & 23807 & 16702 & 16675 & 16701 & 16554 & 16497 & \textbf{16462} & 16502\\
\texttt{BitcoinAlpha+} & 3683 & 12972 & 14223 & 13366 & 16006 & 11625 & 11608 & 11619 & 11503 & 11461 & \textbf{11437} & 11464\\
\texttt{com-Amazon} & 334863 & 925872 & 904427 & 897069 & 1138129 & 632915 & 629054 & 626887 & 617758 & 615008 & \textbf{612313} & 615002\\
\texttt{com-DBLP} & 317080 & 1049866 & 948075 & 933506 & 1166731 & 617025 & 615183 & 614162 & 608743 & 608173 & \textbf{606752} & 608212\\
\texttt{com-Youtube} & 1134890 & 2987624 & 3638780 & 3389320 & -- & 2683459 & \textbf{2682697} & -- & -- & -- & -- & --\\
\bottomrule
\end{tabular}
}
\end{table}

\begin{table}[!tbp]
\centering
\caption{SNAP graphs: lower bounds and certified optimality gaps.  The gap is
$(\mathrm{UB}-\LB)/\LB$, where UB is the best value found in any of our
runs (including the head-to-head runs of Section~\ref{sec:h2h}) and LB is the
CertiFlip bound accepted by the independent checker (one 600\,s run on a
Colab machine, Section~\ref{sec:h2h}).}
\label{tab:snap-bounds}
\small
\begin{tabular}{lrrrrr}
\toprule
Graph & best UB & greedy packing & triangle packing & CertiFlip LB & certified gap (\%)\\
\midrule
\texttt{ca-GrQc} & 6050 & 4479 & 4725 & 6035 & 0.25\\
\texttt{ca-HepTh} & 15813 & 10835 & 10887 & 15699 & 0.73\\
\texttt{ca-HepPh} & 46133 & 35012 & 39495 & 43314 & 6.51\\
\texttt{ca-AstroPh} & 125541 & 85348 & 86009 & 117966 & 6.42\\
\texttt{ca-CondMat} & 56834 & 39380 & 39704 & 55158 & 3.04\\
\texttt{email-Enron} & 150419 & 85082 & 83057 & 140551 & 7.02\\
\texttt{loc-Brightkite} & 174695 & 102317 & 100659 & 163166 & 7.07\\
\texttt{soc-Epinions} & 365137 & -- & -- & -- & --\\
\texttt{Slashdot+} & 351856 & -- & -- & -- & --\\
\texttt{Epinions+} & 521786 & -- & -- & -- & --\\
\texttt{BitcoinOTC+} & 16462 & 9137 & 8774 & 15583 & 5.64\\
\texttt{BitcoinAlpha+} & 11437 & 6335 & 6134 & 10722 & 6.67\\
\texttt{com-Amazon} & 612307 & 412040 & 426755 & 556146 & 10.10\\
\texttt{com-DBLP} & 606723 & 399183 & 426206 & 581805 & 4.28\\
\texttt{com-Youtube} & 2649253 & -- & -- & -- & --\\
\bottomrule
\end{tabular}

\end{table}

\paragraph{Primal quality.}  KaPoCE finds the best solution on every graph.
Our implementation of the iterated flipping local search stays within
$0.10\%$--$0.44\%$ of it.  It takes between one second and three minutes,
depending on the graph, compared with KaPoCE's ten minutes.  Louvain- and
Leiden-type local search, which optimise the same objective, are $0.6\%$ to
$2.7\%$ worse.  Pivot, the best of 50 Pivot runs and MatchFlipPivot are far
from competitive, which is consistent with the PACE results.  CertiFlip
spends half of its budget on the bound, and its gap-guided LNS improves the
flipping solution only slightly on these graphs.  Its primal quality is
therefore that of IteratedFlip.  On four graphs CertiFlip exceeded the budget
by up to a factor of two, because single local-search rounds or sub-MIP
constructions are not interrupted.  This does not change the comparison,
since its solutions there are still worse than KaPoCE's.

\paragraph{Certificates.}  Table~\ref{tab:snap-bounds} is the main result of
this section.  For these bounds CertiFlip exported its certificate, and every
bound in the table is the value computed by the independent checker of
Section~\ref{sec:lean}, which accepted all eleven certificates (the check took
1 to 34\,s).  On the eleven graphs whose support fits into memory, the
CertiFlip bound certifies the best known solutions within $0.25\%$ to $10.1\%$
of optimal.  The median gap is $6.4\%$, and on the collaboration networks
\texttt{ca-GrQc} and \texttt{ca-HepTh} it is below $1\%$.  The classical
scalable bounds, a maximal packing of bad triangles (the bound of
MatchFlipPivot) and the fractional triangle packing, certify only factors of
$1.17$ to $1.81$ on the same graphs.  As far as we know, no certificate of this
quality was previously available for correlation clustering instances of
this size.  The certificates for \texttt{com-Amazon} and \texttt{com-DBLP}
($3\cdot10^5$ vertices, $10^6$ edges) were computed within the 600\,s budget
on one core.  The dense-matrix implementation of the KaPoCE packing bounds cannot
represent these graphs.



\subsection{Ablations}

\paragraph{Components of the bound.}  Table~\ref{tab:ablation} evaluates the
components of the lower bound on five graphs whose size allows running the
dense-matrix KaPoCE bounds for comparison.  Each variant had a budget of
$300$\,s.  Gaps are measured against the best known solution.
\begin{itemize}
\item \emph{Triangle packings.}  Packings of bad triangles, greedy or
fractional, leave gaps of $28\%$ to $88\%$.  This is the factor-two regime of
previous scalable bounds.
\item \emph{Star packings.}  A greedy star packing on the sparse support
already reaches gaps of $5\%$ to $8\%$ in well under a second.  The
ruin-and-recreate local search lowers them to $0.7\%$ to $6.3\%$.
\item \emph{Block LPs.}  Without a warm start, block LPs converge slowly and
reach only an $18\%$ to $75\%$ gap within the budget; \texttt{ca-GrQc} is
the exception.  Started from the packing, they improve the bound in every
case, to $0.2\%$ to $6.2\%$.  Gap-guided blocks add a little more.
\item \emph{Comparison with KaPoCE.}  The root star-packing bound of the
KaPoCE branch-and-bound is weaker than ours on \texttt{ca-GrQc}, equal on
\texttt{ca-HepTh}, and stronger on \texttt{BitcoinAlpha+} ($3.6\%$ against
$6.1\%$), where it needed $850$\,s.  On \texttt{BitcoinOTC+} and
\texttt{ca-CondMat} it did not finish within one hour.
\end{itemize}

\begin{table}[t]
\centering
\caption{Ablation of the lower bound (budget $300$\,s per variant): certified
gap to the best known solution and running time.  ``--'': did not finish
within one hour.}
\label{tab:ablation}
\scriptsize
\setlength{\tabcolsep}{3pt}
\resizebox{\textwidth}{!}{\begin{tabular}{lrrrrrrrrrr}
\toprule
 & \multicolumn{2}{c}{\texttt{ca-GrQc}} & \multicolumn{2}{c}{\texttt{ca-HepTh}} & \multicolumn{2}{c}{\texttt{BitcoinAlpha+}} & \multicolumn{2}{c}{\texttt{BitcoinOTC+}} & \multicolumn{2}{c}{\texttt{ca-CondMat}}\\
Bound & gap (\%) & s & gap (\%) & s & gap (\%) & s & gap (\%) & s & gap (\%) & s\\
\midrule
greedy triangle packing & 35.07 & 0 & 45.95 & 0 & 80.54 & 0 & 80.17 & 0 & 44.32 & 0\\
triangle packing (MWU) & 28.04 & 0 & 45.26 & 0 & 86.45 & 1 & 87.62 & 2 & 43.15 & 3\\
greedy star packing & 5.09 & 3 & 5.96 & 0 & 7.98 & 0 & 7.33 & 0 & 5.58 & 0\\
star packing + local search & 0.73 & 62 & 0.93 & 46 & 6.34 & 39 & 5.62 & 55 & 3.23 & 20\\
block LP (cold start) & 0.78 & 230 & 17.85 & 382 & 47.86 & 361 & 74.94 & 317 & 59.20 & 678\\
packing + block LP & 0.22 & 130 & 0.76 & 304 & 6.16 & 303 & 5.45 & 468 & 2.04 & 666\\
packing + block LP + gap blocks & 0.18 & 300 & 0.75 & 343 & 6.14 & 305 & 5.46 & 470 & 2.03 & 667\\
KaPoCE root p3 packing & 23.19 & 47 & 38.14 & 374 & 76.72 & 9 & -- & -- & -- & --\\
KaPoCE root star packing & 0.58 & 31 & 0.92 & 277 & 3.61 & 850 & -- & -- & -- & --\\
\bottomrule
\end{tabular}
}
\end{table}

\paragraph{Scaling.}  Table~\ref{tab:scaling} reports running times on
synthetic sparse instances with planted clusters of mean size $10$,
intra-cluster density $0.7$ and two random inter-cluster edges per vertex.
Building the distance-two support and one pass of Louvain-type local search
take seconds even at $10^6$ vertices.  The iterated flipping local search
grows almost linearly.  One million ruin-and-recreate steps of the packing
local search take less than a minute beyond the linear-time initial
construction.

\begin{table}[t]
\centering
\caption{Running times (s) on synthetic sparse instances.}
\label{tab:scaling}
\small
\begin{tabular}{lrrrr}
\toprule
Step & $n=10\,000$ & $n=30\,000$ & $n=100\,000$ & $n=300\,000$\\
\midrule
iterated flip & 6.3 & 17.7 & 65.2 & 259.6\\
louvain & 0.1 & 0.1 & 0.4 & 1.6\\
pivot & 0.5 & 0.0 & 0.0 & 0.0\\
star packing (1e6 iterations) & 20.6 & 23.4 & 28.9 & 45.8\\
support & 0.1 & 0.2 & 0.6 & 2.7\\
\bottomrule
\end{tabular}

\end{table}


\subsection{PXMem against KaPoCE}\label{sec:h2h}

\paragraph{Setup.}  The head-to-head runs used Google Colab machines with one
physical core (two hyperthreads) and 12\,GB of memory.  On every machine the
two solvers ran one after the other, each pinned to the same logical CPU while
the rest of the machine was idle, with the same wall-clock budget of
\SI{600}{s}.  KaPoCE was stopped as in PACE~2021, by SIGTERM at the limit; it
normally stops itself after 590\,s.  Its seed, fixed in the released code, was
set to the run seed (a one-line change), so that its ten runs per graph are
independent.  Its output was checked: an edit set is valid iff the edited graph
is a disjoint union of cliques.  A run in which KaPoCE returned no valid
solution would count as a win for PXMem, as in the PACE rules; this did not
occur in any of the 150 runs.  PXMem ran with the released configuration.  The
Numba kernels were compiled once per machine before the timed runs, as KaPoCE
is compiled ahead of time.  We used all fifteen SNAP graphs, ten seeds each.
The configuration of PXMem was developed on eleven of these graphs; the PACE
heuristic-track instances below were not used during development.  We report
the mean paired difference $\bar\Delta$ (PXMem minus KaPoCE, in
disagreements) with a 95\% bootstrap interval, wins, ties and losses, the
two-sided sign test and its Holm correction over the fifteen graphs.

\begin{table}[t]
\centering
\caption{PXMem against KaPoCE, \SI{600}{s} per run, run one after the other on
the same machine and CPU (ten seeds per graph).  Mean costs; $\bar\Delta$: mean
paired difference PXMem$-$KaPoCE with a 95\% bootstrap interval; W/T/L: wins,
ties and losses of PXMem; $p$: two-sided sign test; $p_{\mathrm{Holm}}$: Holm
correction over the fifteen graphs (bold if below $0.05$); $t_{=}$: median time
at which PXMem first reached the final cost of KaPoCE (``--'' if it did not in
most runs).  KaPoCE returned a valid solution in every run.}\label{tab:h2h}
\small
\begin{tabular}{lrrrrrrrr}
\toprule
graph & $n$ & PXMem & KaPoCE & $\bar\Delta$ [95\% CI] & W/T/L & $p$ & $p_{\mathrm{Holm}}$ & $t_{=}$ (s) \\
\midrule
\texttt{ca-GrQc} & \num{5242} & \num{6050.0} & \num{6050.0} & $+0.0$  & 0/10/0 & 1.00 & 1.00 & 69 \\
Bitcoin-Alpha$^+$ & \num{3683} & \num{11437.0} & \num{11437.0} & $+0.0$  & 0/10/0 & 1.00 & 1.00 & 71 \\
Bitcoin-OTC$^+$ & \num{5573} & \num{16462.0} & \num{16462.0} & $+0.0$  & 0/10/0 & 1.00 & 1.00 & 181 \\
\texttt{ca-HepTh} & \num{9877} & \num{15813.0} & \num{15813.5} & $-0.5$ $[-0.8, -0.2]$ & 5/5/0 & 0.06 & 0.62 & 125 \\
\texttt{ca-HepPh} & \num{12008} & \num{46133.0} & \num{46133.0} & $+0.0$  & 0/10/0 & 1.00 & 1.00 & 107 \\
\texttt{ca-AstroPh} & \num{18772} & \num{125543.9} & \num{125606.8} & $-62.9$ $[-78.1, -49.5]$ & 10/0/0 & \textbf{0.002} & \textbf{0.029} & 62 \\
\texttt{ca-CondMat} & \num{23133} & \num{56834.3} & \num{56845.6} & $-11.3$ $[-16.4, -6.0]$ & 8/2/0 & \textbf{0.008} & 0.10 & 56 \\
\texttt{email-Enron} & \num{36692} & \num{150428.2} & \num{150460.8} & $-32.6$ $[-47.5, -17.1]$ & 9/0/1 & \textbf{0.021} & 0.26 & 139 \\
\texttt{loc-Brightkite} & \num{58228} & \num{174700.1} & \num{174703.2} & $-3.1$ $[-6.0, -0.4]$ & 7/0/3 & 0.34 & 1.00 & 432 \\
\texttt{Slashdot+} & \num{75144} & \num{351865.0} & \num{351869.1} & $-4.1$ $[-8.7, +0.1]$ & 6/2/2 & 0.29 & 1.00 & 579 \\
\texttt{soc-Epinions} & \num{75879} & \num{365164.5} & \num{365167.6} & $-3.1$ $[-22.7, +16.1]$ & 4/0/6 & 0.75 & 1.00 & -- \\
\texttt{Epinions+} & \num{114467} & \num{521817.3} & \num{521818.1} & $-0.8$ $[-17.6, +16.6]$ & 5/0/5 & 1.00 & 1.00 & -- \\
\texttt{com-DBLP} & \num{317080} & \num{606753.1} & \num{606787.2} & $-34.1$ $[-56.5, -15.1]$ & 9/0/1 & \textbf{0.021} & 0.26 & 296 \\
\texttt{com-Amazon} & \num{334863} & \num{612328.7} & \num{612311.6} & $+17.1$ $[+14.5, +20.0]$ & 0/0/10 & \textbf{0.002} & \textbf{0.029} & -- \\
\texttt{com-Youtube} & \num{1134890} & \num{2650090.2} & \num{2649791.7} & $+298.5$ $[+40.0, +574.9]$ & 4/0/6 & 0.75 & 1.00 & -- \\
\bottomrule
\end{tabular}

\end{table}

\paragraph{Solution quality at \SI{600}{s}.}  Table~\ref{tab:h2h} shows that on
the four smallest graphs and on \texttt{ca-HepPh} both solvers find the same
solutions in every run; PXMem reaches them after one to three minutes.  PXMem
is better in every run on \texttt{ca-AstroPh} (by 63 disagreements on average)
and in eight to nine of ten runs on \texttt{ca-CondMat}, \texttt{email-Enron}
and \texttt{com-DBLP} (by 11 to 34).  These four differences are significant by
the sign test, but after the Holm correction over fifteen graphs only
\texttt{ca-AstroPh} remains so.  On the four social networks of medium size the
two solvers are statistically indistinguishable.  KaPoCE is better in every run
on \texttt{com-Amazon} (by 17 disagreements, $0.003\%$), which survives the
correction, and on average on \texttt{com-Youtube}, where the runs of both
solvers vary by several hundred disagreements and the sign test is not
significant.  Over all 150 runs PXMem wins 67, ties 49 and loses 34.  The
differences are small in relative terms (at most $0.05\%$), far below the
certified gaps of Section~\ref{sec:snap}; they matter only when the last
fraction of a percent is wanted.  On \texttt{com-Amazon} we compared the two solutions of a local \SI{300}{s}
run block by block with Theorem~\ref{thm:px}.  They differed in more than $16{,}000$ blocks, almost all
ties, and KaPoCE was better in $24$ blocks of $8$ to $77$ vertices, each
containing a vertex of degree $22$ to $142$.  Solving these blocks exactly
with the MIP of Section~\ref{sec:gap} recovered the value of KaPoCE in $23$ of
them, mostly in less than a second, so the difficulty lies in finding the blocks,
not in solving them.

\begin{table}[t]
\centering
\caption{Mean relative difference (\%) of PXMem against KaPoCE for different
budgets (two runs per graph at \SI{60}{s} and \SI{150}{s}, ten at \SI{600}{s});
negative values favour PXMem.}\label{tab:budget}
\small
\begin{tabular}{lrrrrrr}
\toprule
 & \multicolumn{2}{c}{\SI{60}{s}} & \multicolumn{2}{c}{\SI{150}{s}} & \multicolumn{2}{c}{\SI{600}{s}} \\
\cmidrule(lr){2-3}\cmidrule(lr){4-5}\cmidrule(lr){6-7}
graph & med.\ (\%) & W/T/L & med.\ (\%) & W/T/L & med.\ (\%) & W/T/L \\
\midrule
\texttt{ca-GrQc} & $+0.000$ & 0/10/0 & $+0.000$ & 0/10/0 & $+0.000$ & 0/10/0 \\
Bitcoin-Alpha$^+$ & $+0.000$ & 2/6/2 & $+0.000$ & 0/8/2 & $+0.000$ & 0/10/0 \\
Bitcoin-OTC$^+$ & $+0.009$ & 0/2/8 & $+0.000$ & 0/9/1 & $+0.000$ & 0/10/0 \\
\texttt{ca-HepTh} & $-0.006$ & 7/2/1 & $-0.006$ & 8/2/0 & $-0.003$ & 5/5/0 \\
\texttt{ca-HepPh} & $-0.002$ & 8/2/0 & $+0.000$ & 2/8/0 & $+0.000$ & 0/10/0 \\
\texttt{ca-AstroPh} & $+0.037$ & 3/0/7 & $+0.009$ & 4/0/6 & $-0.047$ & 10/0/0 \\
\texttt{ca-CondMat} & $-0.019$ & 7/1/2 & $-0.030$ & 9/0/1 & $-0.029$ & 8/2/0 \\
\texttt{email-Enron} & $-0.001$ & 5/0/5 & $-0.013$ & 7/0/3 & $-0.025$ & 9/0/1 \\
\texttt{loc-Brightkite} & $-0.001$ & 5/0/5 & $-0.001$ & 5/1/4 & $-0.001$ & 7/0/3 \\
\texttt{Slashdot+} & $+0.018$ & 0/0/10 & $+0.002$ & 3/0/7 & $-0.001$ & 6/2/2 \\
\texttt{soc-Epinions} & $+0.016$ & 1/0/9 & $+0.004$ & 2/0/8 & $+0.000$ & 4/0/6 \\
\texttt{Epinions+} & $+0.010$ & 1/0/9 & $+0.007$ & 1/0/9 & $+0.000$ & 5/0/5 \\
\texttt{com-DBLP} & $+0.018$ & 1/0/9 & $+0.013$ & 0/0/10 & $-0.004$ & 9/0/1 \\
\texttt{com-Amazon} & $+0.033$ & 1/0/9 & $+0.012$ & 0/0/10 & $+0.003$ & 0/0/10 \\
\texttt{com-Youtube} & $+0.641$ & 2/0/8 & $+0.188$ & 0/0/10 & $+0.009$ & 4/0/6 \\
\bottomrule
\end{tabular}

\end{table}

\paragraph{Dependence on the budget.}  Table~\ref{tab:budget} shows the
limitation of PXMem.  On the two largest graphs KaPoCE is $0.1\%$ to $0.7\%$
better at \SI{60}{s} and \SI{150}{s}, because the initial population (four
Pivot, flipping and annealing runs) takes a large part of a short budget.  On
\texttt{ca-HepTh}, \texttt{ca-AstroPh} and \texttt{ca-CondMat} PXMem is better at
every budget.

\paragraph{Parallel runs and post-processing.}  Theorem~\ref{thm:px} gives a
way to use several cores that heuristics without an exact recombination do not
have.  Table~\ref{tab:parallel} reports four independent PXMem runs of
\SI{300}{s} that ran at the same time on four cores of our 4-core machine,
compared with the best of $k$ of them and with their partition crossover, and
with KaPoCE on one core.  Combining two runs by partition crossover is better
than the best of four runs on four of the five graphs where the runs differ
(all but \texttt{email-Enron}), and never worse than the best of two.  With four
runs PXMem matches or beats KaPoCE on five of the six graphs and is one
disagreement behind on \texttt{com-Amazon}.  Applied to the solutions of
KaPoCE, the crossover improves them on all five graphs where they are not
already equal to ours (by $2$ to $51$ disagreements), which makes it useful as
a post-processing step for any solver.  These numbers compare four cores with
one and complement, but do not replace, the single-core comparison of
Table~\ref{tab:h2h}.

\begin{table}[t]
\centering
\caption{Four independent PXMem runs (\SI{300}{s} each, run in parallel):
best of $k$ runs and partition crossover of $k$ runs (averaged over all
subsets of size $k$), KaPoCE (\SI{300}{s}, one core), and the partition
crossover of the KaPoCE solution with a PXMem solution (averaged over the four
runs).  The best value in each row is bold.}\label{tab:parallel}
\small
\begin{tabular}{lrrrrrrrr}
\toprule
 & 1 run & \multicolumn{2}{c}{2 runs} & \multicolumn{2}{c}{4 runs} & & \multicolumn{2}{c}{PX with KaPoCE} \\
\cmidrule(lr){3-4}\cmidrule(lr){5-6}\cmidrule(lr){8-9}
graph & & best & PX & best & PX & KaPoCE & 1 run & 4 runs \\
\midrule
Bitcoin-OTC$^+$ & \textbf{\num{16462}} & \textbf{\num{16462}} & \textbf{\num{16462}} & \textbf{\num{16462}} & \textbf{\num{16462}} & \textbf{\num{16462}} & \textbf{\num{16462}} & \textbf{\num{16462}} \\
\texttt{ca-AstroPh} & \num{125544.5} & \num{125542.8} & \num{125541.7} & \num{125542} & \textbf{\num{125540}} & \num{125567} & \num{125542.8} & \textbf{\num{125540}} \\
\texttt{email-Enron} & \num{150453.8} & \num{150445.8} & \num{150444.5} & \num{150433} & \textbf{\num{150432}} & \num{150438} & \num{150435.8} & \textbf{\num{150432}} \\
\texttt{loc-Brightkite} & \num{174706.2} & \num{174704.5} & \num{174699.7} & \num{174703} & \num{174697} & \num{174699} & \num{174696} & \textbf{\num{174695}} \\
\texttt{com-DBLP} & \num{606773.8} & \num{606761.2} & \num{606735.3} & \num{606754} & \num{606720} & \num{606779} & \num{606728.2} & \textbf{\num{606719}} \\
\texttt{com-Amazon} & \num{612374.5} & \num{612361.7} & \num{612330.3} & \num{612350} & \num{612317} & \num{612316} & \num{612309.2} & \textbf{\num{612308}} \\
\bottomrule
\end{tabular}

\end{table}

\paragraph{Memory.}  PXMem needs more memory than KaPoCE.  The peak resident
set sizes measured in the runs of Table~\ref{tab:h2h} are \SI{360}{MB} to
\SI{370}{MB} against \SI{180}{MB} to \SI{190}{MB} on \texttt{com-Amazon} and
\texttt{com-DBLP}, and \SI{970}{MB} against \SI{730}{MB} on \texttt{com-Youtube}.



\section{Conclusion}

Correlation clustering has moved quickly in theory, but certified solution
quality on real instances had not kept pace.  Three ingredients make LP-based
certificates usable at scale: the distance-two relaxation, the anytime
block-dual bound warm-started by sparse packings, and the gap decomposition.
They also give a primal method that knows where to search.  Our
implementation of the iterated flipping local search of
\citet{CohenAddadLPTYZ24} is, to our knowledge, the first of a sub-$3$
combinatorial scheme.  It comes within a few tenths of a percent of the
best engineered heuristic in a fraction of its running time.  The exact
partition crossover closes the remaining gap on most graphs: PXMem keeps the
a-priori guarantee, is significantly better than KaPoCE on four of eleven
graphs at equal budget, and turns additional cores into better solutions.  All
combinatorial arguments are machine-checked.

\paragraph{Limitations.}
\begin{itemize}
\item KaPoCE remains better than PXMem on \texttt{com-Amazon} (by $0.004\%$ at
\SI{600}{s}) and, at short budgets of \SI{60}{s}--\SI{150}{s}, on the largest
graphs (by $0.1\%$--$0.7\%$), where our initial population takes too much of
the budget.  PXMem also needs about $2.5$ times the memory of KaPoCE.
\item The Pivot guarantee used in Theorems~\ref{thm:certiflip}
and~\ref{thm:pxmem} is taken from the literature; our Lean development proves
everything else in these theorems but not this probabilistic argument.
\item On dense instances the packing bounds of the KaPoCE branch-and-bound
remain stronger than ours.
\item The distance-two support has $O(\sum_v\deg(v)^2)$ pairs.  On graphs
with strong hubs it exceeds the memory of our machine, and CertiFlip then
reports no bound.  Restricting the support to centres of bounded degree keeps
the bound valid and is a natural next step.
\item The $2.06$ a-priori guarantee applies only when a single block covers
the graph and separation converges.
\end{itemize}

\paragraph{Open directions.}  Parallel block solves; cluster-LP columns
inside blocks, where the recent roundings of \citet{CaoCLLNV24} and
\citet{GarciaSorianoS26} could be applied; and overlapping blocks with
Lagrangian cost splitting.

\paragraph{Code and data availability.}  All code, instance lists and raw
results are available at
\url{https://github.com/vinhqdang/correlation_clustering}.

\bibliographystyle{plainnat}
\bibliography{refs}

\end{document}
